\section*{Bab \ref{ch:b1:radiometric-methods} --- Metode penanggalan radiometrik}
% nilai: calc/ch25.py

\begin{solution}{exo:b1:radiometric-methods:1}
(a)~$D^*/P = 2^n - 1 = 3$ memberi $n = 2$: dua waktu paruh. (b)~Rubidium-87:
$e^{\lambda t} - 1 = e^{\num{1.3972e-11} \times \num{3.10e8}} - 1 = 0.004341$;
uranium-235: $e^{\num{9.8485e-10} \times \num{3.10e8}} - 1 = 0.3570$. Rasio uranium
sekitar 80 kali lebih besar: galat mutlak pengukuran tertentu adalah bagian yang jauh
lebih kecil darinya. (c)~$2^{-n} < 0.01$ ketika $n > \log_2 100 = 6.64$: setelah
sekitar tujuh waktu paruh.
\end{solution}

\begin{solution}{exo:b1:radiometric-methods:2}
(a)~\omterm{def:b1:radiometric-methods:radiocarbon}{Radiokarbon} pada arang. (b)~Uranium--timbal pada zirkon, atau argon--argon
(kalium--argon) pada sanidin, \omterm{def:b1:rock-forming-minerals:mineral}{mineral} kaya kalium. (c)~Uranium--timbal pada zirkon, yang
bertahan dari sejarah panjang; rubidium--stronsium pada seluruh-batuan mungkin; kalium--argon
berisiko memberi umur yang turun karena argon hilang dalam pemanasan kemudian. (d)~\omterm{def:b1:radiometric-methods:radiocarbon}{Radiokarbon}
pada karbonat cangkang, dekat batas metode (sekitar \num{70000} tahun). (e)~Tidak ada
secara langsung: umurnya diapit oleh batuan vulkanik bertarikh di atas dan di bawah, atau
ditemukan dengan mengkorelasikan fosilnya dengan penampang bertarikh (\cref{ch:b1:stratigraphic-scale}).
\end{solution}

\begin{solution}{exo:b1:radiometric-methods:3}
(a)~$t = -8033 \ln 0.25 = 8033 \times 1.3863 = \num{11136}$ tahun BP.
(b)~Umur rata-ratanya adalah $5700/\ln 2 = \num{8223}$ tahun: $t = 8223 \times 1.3863 =
\num{11400}$ tahun. (c)~Keduanya mengasumsikan bahwa proporsi karbon-14 dalam karbon hidup
pada waktu itu sama dengan standar; itu tidak. Umur kalender dibaca dari \omterm{def:b1:radiometric-methods:radiocarbon}{kurva kalibrasi}
seperti IntCal20, sebagai rentang tahun cal BP.
\end{solution}

\begin{solution}{exo:b1:radiometric-methods:4}
(a)~Setiap jam mulai ketika \omterm{def:b1:rock-forming-minerals:mineral}{mineralnya} mendingin di bawah \omterm{def:b1:radiometric-methods:closure}{suhu penutupannya}: zirkon (di
atas \qty{750}{\celsius}) menutup lebih dulu, hornblenda (sekitar \qty{530}{\celsius})
berikutnya, biotit (sekitar \qty{280}{\celsius}) terakhir, sehingga umurnya menurun dalam
urutan itu seiring granit mendingin. (b)~$(530 - 280)/(410 - 398) =
250/12 = \qty{21}{\celsius}$ per juta tahun. (c)~Zirkon menutup di atas
\qty{750}{\celsius}, sehingga granit mendingin sekurang-kurangnya $750 - 530 =
\qty{220}{\celsius}$ dalam \qty{10}{\Ma}: sekurang-kurangnya \qty{22}{\celsius} per juta
tahun, sebuah batas bawah.
\end{solution}

\begin{solution}{exo:b1:radiometric-methods:5}
(a)~$\bar x = 7.5$, $\bar y = 0.809525$, $S_{xx} = 81.0$, $S_{xy} = 1.13965$;
$m = 1.13965/81.0 = 0.0140698$, $b = 0.809525 - 7.5 \times 0.0140698 = 0.70400$.
(b)~$t = \ln(1.0140698)/\num{1.3972e-11} = \qty{1000.0}{\Ma}$; \omterm{def:b1:radiometric-methods:isochron}{rasio awal}
$0.7040$. (c)~Sisaannya adalah $+1.6$, $-2.5$, $+0.7$ dan $+0.2$ (dalam satuan
\num{e-4}): $s = \sqrt{\num{9.34e-8}/2} = 0.00022$, $\sigma_m = 0.00022/\sqrt{81}
= \num{2.40e-5}$ dan $\sigma_t = \num{2.40e-5}/(\num{1.3972e-11} \times 1.0141) =
\qty{1.7}{\Ma}$: \qty[uncertainty-mode=separate]{1000.0(17)}{\Ma}. (d)~$0.7040$
dekat dengan nilai \omterm{def:b1:layered-earth:layers}{mantel} $0.7025$ dan jauh di bawah \omterm{def:b1:layered-earth:layers}{kerak} $0.710$: sumber sebagian besar
di \omterm{def:b1:layered-earth:layers}{mantel}.
\end{solution}

\begin{solution}{exo:b1:radiometric-methods:6}
(a)~$0.100/39.098 \times \num{1.17e-4} = \qty{2.992e-7}{mol}$ kalium-40 per gram.
(b)~${}^{40}\mathrm{Ar}^*/{}^{40}\mathrm{K} = \num{1.169e-9}/\num{2.992e-7}
= 0.0039065$; $\lambda/\lambda_e = 5.543/0.581 = 9.5404$; $t = \ln(1 + 9.5404
\times 0.0039065)/\num{5.543e-10} = \ln(1.037269)/\num{5.543e-10} =
\qty{66.01}{\Ma}$. (c)~Argon radiogenik akan menjadi $0.995 \times 1.169$
nanomol, memberi \qty{65.69}{\Ma}: umur yang diukur akan \qty{0.32}{\Ma} terlalu tua,
sepuluh kali jendela \num{32000} tahun tempat tumbukan dan kepunahan di akhir Kapur kini
ditunjukkan berimpit. Argon berlebih menjadi masalah terutama pada batuan muda dan
seluruh-batuan; \omterm{def:b1:rock-forming-minerals:mineral}{mineral} seperti sanidin seharusnya tidak memerangkap argon ketika ia
mengkristal, dan spektrum umur argon--argon akan mengungkapkannya.
\end{solution}

\begin{solution}{exo:b1:radiometric-methods:7}
(a)~Sekitar \num{2000} cal BP: kurva melewati \num{2070} tahun BP di sana.
(b)~Dari sekitar \num{2495} hingga \num{2699} cal BP, dua ratus tahun: kurva hampir datar
di antara \num{2500} dan \num{2700} cal BP, sehingga setiap umur kalender dalam interval
itu memiliki \omterm{def:b1:radiometric-methods:radiocarbon}{umur radiokarbon} yang hampir sama, dan pengukuran yang lebih tepat tidak akan
mempersempit rentangnya. (c)~Pada \num{2000} cal BP \omterm{def:b1:radiometric-methods:radiocarbon}{umur radiokarbon} adalah \num{70} tahun
\emph{lebih tua} daripada umur kalender; pada \num{8000} cal BP \num{822} tahun lebih muda
(\qty{10}{\percent}); pada \num{12000} cal BP \num{1694} tahun lebih muda
(\qty{14}{\percent}). Perbedaannya tidak sebanding dengan umur: ia bahkan berubah tanda. Ia
tidak dapat berasal dari waktu paruh saja, yang akan memberi \qty{2.4}{\percent} yang
konstan.
\end{solution}

\begin{solution}{exo:b1:radiometric-methods:8}
(a)~Lapisan abu di bawah (\qty{94.50}{\Ma}) memberi umur maksimum, aliran lava di atas
(\qty{93.80}{\Ma}) umur minimum: batas itu diapit di antara keduanya. Dike juga umur
minimum, tetapi yang lebih lemah, lebih muda daripada lava. Zirkon detrital hanya umur
maksimum bagi batu pasir, dan lebih lemah daripada lapisan abu, yang terletak di bawah
batu pasir. (b)~$f = (12 - 4)/(20 - 4) = 0.5$; $t = 94.50 - 0.5 \times 0.70 =
\qty{94.15}{\Ma}$; $\sigma_t = \sqrt{0.25 \times 0.10^2 + 0.25 \times 0.12^2} =
\qty{0.08}{\Ma}$: \qty[uncertainty-mode=separate]{94.15(8)}{\Ma}. (c)~Butirannya berasal
dari batuan yang berumur sekurang-kurangnya \qty{95.3}{\Ma}, tererosi di suatu tempat di
daerah sumber; ini menunjukkan bahwa batu pasir lebih muda daripada \qty{95.3}{\Ma}, yang
sudah dibuktikan lebih ketat oleh lapisan abu di bawahnya.
\end{solution}

\begin{solution}{exo:b1:radiometric-methods:9}
(a)~Butir~A: $\ln(1.08401)/\num{1.55125e-10} = \qty{520.0}{\Ma}$ dan
$\ln(1.6688)/\num{9.8485e-10} = \qty{520.0}{\Ma}$. Butir~B: \qty{368.4}{\Ma} dan
\qty{389.9}{\Ma}. (b)~Butir~A konkordan. (c)~$0.05881/0.08401 = 0.7000$ dan
$0.4682/0.6688 = 0.7001$. Umurnya adalah $\ln(1 + R)/\lambda$, tidak sebanding dengan
$R$: untuk rasio uranium-238 yang kecil $\ln(1+R) \approx R$ dan umurnya turun hampir
\qty{30}{\percent} (menjadi \qty{368}{\Ma}); untuk rasio uranium-235 yang lebih besar
logaritmanya turun lebih sedikit, dan umurnya hanya menjadi \qty{390}{\Ma}. Keduanya
terlalu muda, dan keduanya tidak sepakat: butir itu tidak konkordan.
\end{solution}

\begin{solution}{exo:b1:radiometric-methods:10}
(a)~Dengan $u = \lambda t$, $D^*/P = e^u - 1 = u + u^2/2 + \dots$, sehingga $(D^*/P)/\lambda
= t\,(1 + u/2 + \dots)$: pendekatan itu melebih-lebihkan umur dengan galat relatif yang
mendekati $u/2$. (b)~Rubidium-87 pada \qty{310}{\Ma}: $u = 0.00433$, galat
\qty{0.22}{\percent} (\qty{0.7}{\Ma}); kalium-40 pada \qty{296}{\Ma}: $u =
0.164$, galat \qty{8.7}{\percent} (\qty{26}{\Ma}; $u/2$ memberi \qty{8.2}{\percent});
uranium-238 pada \qty{4.4}{\Ga}: $u = 0.683$, galat \qty{43}{\percent} --- pendekatan dan
bahkan perkiraan galatnya ($u/2 = \qty{34}{\percent}$) gagal. (c)~$R = 0.004341$, $\sigma_R = \num{4.34e-6}$; dengan
\cref{prop:b1:radiometric-methods:precision}, $\sigma_t =
\num{4.34e-6}/(\num{1.3972e-11} \times 1.0043) = \qty{0.31}{\Ma}$.
\end{solution}

\begin{solution}{exo:b1:radiometric-methods:11}
(a)~Misalkan ${}^{40}\mathrm{K}$ adalah kalium-40 hari ini. Pada waktu $t_2$ sebelum
sekarang ada ${}^{40}\mathrm{K}\,e^{\lambda t_2}$ atom, dan argon yang dibuat antara
$t_1$ dan $t_2$ adalah, dengan \cref{prop:b1:radiometric-methods:kar} yang diterapkan pada
interval itu, $(\lambda_e/\lambda)\,{}^{40}\mathrm{K}\,e^{\lambda
t_2}\bigl(e^{\lambda (t_1 - t_2)} - 1\bigr) = (\lambda_e/\lambda)\,
{}^{40}\mathrm{K}\,\bigl(e^{\lambda t_1} - e^{\lambda t_2}\bigr)$. Fraksi
$1 - f$ darinya tetap. Sejak $t_2$, \omterm{def:b1:rock-forming-minerals:mineral}{mineral} yang tertutup telah membuat
$(\lambda_e/\lambda)\,{}^{40}\mathrm{K}\,(e^{\lambda t_2} - 1)$ lagi. Menjumlahkan dan
mengalikan dengan $\lambda/(\lambda_e\,{}^{40}\mathrm{K})$ memberi ungkapan itu.
(b)~Umur semunya adalah $\ln(1 + \text{ungkapan})/\lambda$: \qty{400.0}{\Ma} untuk
$f = 0$, \qty{342.3}{\Ma} untuk $0.3$, \qty{282.7}{\Ma} untuk $0.6$ dan
\qty{200.0}{\Ma} untuk $1$. (c)~Setiap umur antara \qty{200}{\Ma} dan
\qty{400}{\Ma} dapat dihasilkan oleh suatu kehilangan sebagian, dan satu angka tunggal
tidak dapat mengatakan apakah ia mencatat peristiwa nyata pada waktu itu atau campuran;
spektrum umur, atau jam lain dalam batuan yang sama, diperlukan.
\end{solution}

\begin{solution}{exo:b1:radiometric-methods:12}
(a)~Umurnya sebanding dengan umur rata-rata yang dipakai: $10306 \times 8223/8033 =
\num{10550}$ tahun. (b)~$10550 - 10306 = 244$ tahun, sekitar sepertujuh dari
\num{1694}. (c)~Jika aktivitas awalnya adalah $k A_0$, umur sebenarnya adalah $t = \tau
\ln(k A_0/A) = \num{10550} + \tau \ln k$; dengan $t = \num{12000}$, $\ln k =
1450/8223 = 0.176$ dan $k = 1.19$: sekitar \qty{19}{\percent} lebih banyak karbon-14
daripada standar.
\end{solution}

\begin{solution}{pb:b1:radiometric-methods:1}
\begin{enumerate}
  \item Titik-titik naik secara mantap dengan $x$; simpangannya dari suatu garis paling
        banyak dua satuan pada desimal keempat, terhadap rentang $0.04$ dalam $y$: mereka
        terletak pada garis lurus dengan sebaran kecil.
  \item $\bar x = 4.90333$, $\bar y = 0.729288$, $S_{xx} = 62.1253$, $S_{xy} =
        0.269335$; $m = 0.269335/62.1253 = 0.0043354$.
  \item $b = 0.729288 - 0.0043354 \times 4.90333 = 0.70803$: $^{87}\mathrm{Sr}/^{86}\mathrm{Sr}$
        granit ketika ia mengkristal, yang dimiliki bersama oleh semua bagiannya.
  \item $t = \ln(1.0043354)/\num{1.3972e-11} = \qty{309.6}{\Ma}$.
  \item Sisaannya adalah $+9.1$, $-17.2$, $+6.9$, $-6.4$, $+15.1$ dan $-7.5$ (dalam satuan
        \num{e-5}); $s = 0.000137$, $\sigma_m = 0.000137/\sqrt{62.1253} =
        \num{1.74e-5}$, $\sigma_t = \num{1.74e-5}/(\num{1.3972e-11} \times 1.0043)
        = \qty{1.2}{\Ma}$: \qty[uncertainty-mode=separate]{309.6(12)}{\Ma}, dan
        $\sigma_b = 0.00010$.
  \item $0.7080$ lebih dekat ke \omterm{def:b1:layered-earth:layers}{kerak} $0.710$ daripada ke \omterm{def:b1:layered-earth:layers}{mantel} $0.7025$. Untuk campuran,
        $(0.70803 - 0.7025)/(0.710 - 0.7025) = 0.74$: sekitar tiga perempat stronsium berasal
        dari \omterm{def:b1:layered-earth:layers}{kerak} yang meleleh.
  \item $0.0740/39.098 \times \num{1.17e-4} = \qty{2.214e-7}{mol}$ per gram, dikali
        $\num{6.022e23}$: \num{1.33e17} atom.
  \item $^{40}\mathrm{Ar}^*/^{40}\mathrm{K} = \num{4.139e-9}/\num{2.2144e-7} =
        0.018691$; $1 + 9.5404 \times 0.018691 = 1.17832$; $t =
        \ln(1.17832)/\num{5.543e-10} = \qty{296.0}{\Ma}$.
  \item $^{40}\mathrm{K} = \qty{3.292e-8}{mol.g^{-1}}$, rasio $0.019239$,
        $\ln(1.18355)/\num{5.543e-10} = \qty{304.0}{\Ma}$.
  \item $\ln(1.018691)/\num{5.543e-10} = \qty{33.4}{\Ma}$, sembilan kali terlalu muda,
        karena hanya sekitar satu dari sepuluh peluruhan memberi argon. Kalsium-40 adalah
        isotop umum kalsium, melimpah dalam batuan: yang sedikit dibuat oleh peluruhan tidak
        dapat dibedakan dari kalsium yang telah dikandung \omterm{def:b1:rock-forming-minerals:mineral}{mineral} sejak awal.
  \item Argon berlebih menggeser umur sebesar orde satu juta tahun, yang penting untuk batuan
        muda; di sini itu akan jauh di bawah \qty{1}{\percent} dari umurnya. Dan kristal \omterm{def:b1:rock-forming-minerals:groups}{mika}
        serta \omterm{def:b1:rock-forming-minerals:groups}{amfibol} hampir tidak memerangkap argon ketika terbentuk.
  \item Entah granit mendingin perlahan dan \omterm{def:b1:rock-forming-minerals:groups}{mika} serta \omterm{def:b1:rock-forming-minerals:groups}{amfibol} menutup setelah
        kristalisasi (\omterm{def:b1:radiometric-methods:closure}{umur pendinginan}), atau pemanasan kemudian mengeluarkan sebagian argon
        mereka. Plato datar dalam spektrum umur argon--argon mendukung \omterm{def:b1:radiometric-methods:closure}{umur pendinginan};
        spektrum tanpa plato memperingatkan bahwa \omterm{def:b1:rock-forming-minerals:mineral}{mineral-mineral} itu kehilangan argon. Urutan
        umur, hornblenda sebelum biotit seperti yang diramalkan oleh \omterm{def:b1:radiometric-methods:closure}{suhu penutupannya}, juga
        menunjuk pada pendinginan.
  \item $\ln(1.049361)/\num{1.55125e-10} = \qty{310.6}{\Ma}$ dan
        $\ln(1.35784)/\num{9.8485e-10} = \qty{310.6}{\Ma}$: konkordan.
  \item Zirkon \qty{310.6}{\Ma}, penutupan di atas \qty{750}{\celsius}: kristalisasi;
        \omterm{def:b1:radiometric-methods:isochron}{isokron} seluruh-batuan, \qty[uncertainty-mode=separate]{309.6(12)}{\Ma}, juga
        menanggalnya, dan sepakat dalam ketidakpastiannya; hornblenda \qty{304.0}{\Ma}:
        pendinginan melewati sekitar \qty{530}{\celsius}; biotit \qty{296.0}{\Ma}: pendinginan
        melewati sekitar \qty{280}{\celsius}. Grafiknya turun dari lebih dari
        \qty{750}{\celsius} pada \qty{310.6}{\Ma} hingga \qty{280}{\celsius} pada
        \qty{296.0}{\Ma}.
  \item $(530 - 280)/(304.02 - 296.03) = 250/7.99 = \qty{31}{\celsius}$ per juta
        tahun.
  \item $(750 - 530)/(310.60 - 304.02) = 220/6.58 = \qty{33}{\celsius}$ per juta tahun
        setidaknya, karena zirkon menutup di atas \qty{750}{\celsius}: pendinginan setidaknya
        sama cepatnya di awal.
  \item Selisih suhu, $250 \pm \sqrt{40^2 + 40^2} = 250 \pm
        57$, memberi \qtyrange{24}{38}{\celsius} per juta tahun.
  \item $(550 - 350)/7.99 = \qty{25}{\celsius}$ per juta tahun, di dalam rentang soal
        17: ketidakpastiannya berasal dari \omterm{def:b1:radiometric-methods:closure}{suhu penutupan}, bukan dari umur.
  \item $(280 - 10)/30 = \qty{9.0}{km}$.
  \item $\ln(1.039108)/\num{1.55125e-10} = \qty{247.30}{\Ma}$, dan yang sama dari rasio
        uranium-235: konkordan.
  \item Dari granit, yang tererosi pada \omterm{def:b1:cross-sections:unconformity}{ketidakselarasan}: umurnya adalah umur
        kristalisasinya. Keduanya hanya menunjukkan bahwa batu pasir lebih muda daripada
        \qty{310}{\Ma}: umur maksimum, jauh lebih lemah daripada lapisan abu.
  \item $(58 - 18)/(247.30 - 245.30) = \qty{20.0}{m}$ per juta tahun; dasarnya
        $18/20.0 = \qty{0.90}{\Ma}$ lebih tua daripada lapisan abu bawah:
        \qty{248.20}{\Ma}.
  \item Dari \qty{309.6}{\Ma} (\omterm{def:b1:radiometric-methods:isochron}{isokron}) atau \qty{310.6}{\Ma} (zirkon) hingga
        \qty{248.2}{\Ma}: sekitar \qty{62}{\Ma}.
  \item \qty{9.0}{km} dalam $296.0 - 248.2 = \qty{47.8}{\Ma}$: \qty{0.19}{km} per juta
        tahun, \qty{0.19}{mm} per tahun. Ini mengasumsikan bahwa biotit yang ditanggal berasal
        dari tepat di bawah \omterm{def:b1:cross-sections:unconformity}{ketidakselarasan}, sehingga seluruh batuan di atasnya dalam granit
        dan sekis telah tersingkirkan pada saat pasir tiba, dan bahwa \omterm{def:b1:heat-budget:geotherm}{geoterm} tetap konstan.
  \item Granit mendingin sekitar \qty{31}{\celsius} per juta tahun di antara penutupan
        hornblenda dan biotit.
\end{enumerate}
\end{solution}
